\chapter{Wymagania systemowe}
\label{ch:wymagania}

\section{Wymagania funkcjonalne}

Wymagania funkcjonalne okre\'{s}laj\k{a} konkretne zachowania systemu z~perspektywy
ka\.{z}dej z~tr\'{o}jki r\'{o}l u\.{z}ytkownik\'{o}w.

\subsection{Administrator}
\begin{itemize}
  \item Logowanie do systemu przy u\.{z}yciu loginu i~has\l{}a.
  \item Dodawanie, edytowanie i~usuwanie kursant\'{o}w (imi\k{e}, nazwisko, PESEL, login,
        has\l{}o, kategoria prawa jazdy).
  \item Dodawanie, edytowanie i~usuwanie instruktor\'{o}w (imi\k{e}, nazwisko, numer
        legitymacji, telefon, login, has\l{}o, uprawnienia do kategorii).
  \item Dodawanie, edytowanie i~usuwanie pojazd\'{o}w (marka, model, numer rejestracyjny,
        kategoria).
  \item Przegl\k{a}danie pe\l{}nego kalendarza rezerwacji z~mo\.{z}liwo\'{s}ci\k{a}
        filtrowania po kursancie, instruktorze i~poje\'{z}dzie.
  \item Rezerwowanie i~anulowanie jazd w~imieniu dowolnego kursanta.
\end{itemize}

\subsection{Instruktor}
\begin{itemize}
  \item Logowanie do systemu.
  \item Przegl\k{a}danie w\l{}asnego grafiku na bie\.{z}\k{a}cy tydzie\'{n} i~tygodnie
        s\k{a}siednie.
  \item Anulowanie w\l{}asnych zarezerwowanych jazd.
\end{itemize}

\subsection{Kursant}
\begin{itemize}
  \item Logowanie do systemu.
  \item Przegl\k{a}danie kalendarza dost\k{e}pnych termin\'{o}w jazd.
  \item Rezerwowanie wolnego terminu jazdy.
  \item Anulowanie w\l{}asnej zarezerwowanej jazdy.
\end{itemize}

\section{Za\l{}o\.{z}enia dotycz\k{a}ce dzia\l{}ania systemu}

\begin{itemize}
  \item Aplikacja przeznaczona jest do u\.{z}ytku lokalnego na jednym komputerze
        z~systemem Windows.
  \item Ka\.{z}dy u\.{z}ytkownik musi si\k{e} zalogowa\'{c} -- dost\k{e}p bez
        logowania nie jest mo\.{z}liwy.
  \item Dane przechowywane s\k{a} w~lokalnej bazie danych -- nie jest wymagane
        po\l{}\k{a}czenie z~internetem.
  \item Aplikacja obs\l{}uguje jednocze\'{s}nie jedn\k{a} aktywn\k{a} sesj\k{e}
        u\.{z}ytkownika.
  \item Baza danych tworzona jest automatycznie przy pierwszym uruchomieniu.
\end{itemize}

\section{Wymagania niefunkcjonalne}

\begin{table}[H]
\centering
\caption{Wymagania niefunkcjonalne systemu}
\label{tab:wymagania_nf}
\begin{tabularx}{\linewidth}{|l|X|}
\hline
\textbf{Kategoria} & \textbf{Opis} \\
\hline
Bezpiecze\'{n}stwo &
  Has\l{}a u\.{z}ytkownik\'{o}w s\k{a} przechowywane w~zaszyfrowanej formie.
  Loginy musz\k{a} by\'{c} unikalne w~obr\k{e}bie ka\.{z}dej roli (kursanci, instruktorzy, administratorzy osobno). Numery PESEL musz\k{a} by\'{c} unikalne w\'{s}r\'{o}d kursant\'{o}w. \\
\hline
Niezawodno\'{s}\'{c} &
  B\l{}\k{e}dy s\k{a} przechwytywane i~komunikowane u\.{z}ytkownikowi
  w~formie czytelnych komunikat\'{o}w. \\
\hline
Wydajno\'{s}\'{c} &
  Czas odpowiedzi interfejsu poni\.{z}ej~1~sekundy dla typowych operacji. \\
\hline
Przenoszalno\'{s}\'{c} &
  Aplikacja dzia\l{}a na systemie Windows~10 i~Windows~11. \\
\hline
Utrzymywalno\'{s}\'{c} &
  Kod aplikacji jest podzielony na warstwy, co u\l{}atwia rozbudow\k{e}
  o~nowe funkcje. \\
\hline
\end{tabularx}
\end{table}